\chapter{Lattice-Based Digital Signatures}
\label{ch:lattice-signatures}

\abstract*{Constructing digital signatures from lattices requires fundamentally different techniques than encryption. This chapter develops two paradigms: the Fiat-Shamir-with-aborts approach underlying ML-DSA (CRYSTALS-Dilithium) and the hash-and-sign paradigm with Gaussian sampling used by Falcon. Both are NIST standards, and their contrasting designs illuminate deep tradeoffs in lattice-based cryptographic engineering.}

% =============================================================
% SOURCE: notes.tex §4.4.3–4.4.4 (lines 6145–6425)
% REUSE: ~35%. Currently overview-level. Largest expansion in Part II.
% NEW: Fiat-Shamir with aborts, rejection sampling, NTRU lattices,
%      GPV framework, QROM security.
% =============================================================

\section{The Challenge of Lattice-Based Signatures}
\label{sec:sig-challenge}

% TODO: New section
% - Why "encrypt in reverse" doesn't work for lattices
% - The information leakage problem: signatures leak secret key info
% - Two solutions: rejection sampling (Dilithium) and preimage sampling (Falcon)

\section{The Fiat-Shamir Framework}
\label{sec:fiat-shamir}

% TODO: New section
% - Interactive identification protocols
% - The Fiat-Shamir transform: removing interaction with a hash
% - Security in the random oracle model
% - Extension to the quantum random oracle model (QROM)
% - Cite: Fiat-Shamir 1986, Kiltz et al. 2018

\section{Fiat-Shamir with Aborts}
\label{sec:fs-aborts}

% TODO: New section (core technique for Dilithium)
% - Lyubashevsky's framework: commit, challenge, respond, reject
% - Why rejection sampling is necessary: preventing secret leakage
% - The rejection sampling lemma
% - Uniform distribution of signatures regardless of secret key
% - Cite: Lyubashevsky 2009, 2012

\section{ML-DSA (CRYSTALS-Dilithium)}
\label{sec:ml-dsa}

% TODO: Expand from notes.tex §4.4.3
% - KeyGen, Sign, Verify: complete algorithms
% - Module-LWE and Module-SIS foundations
% - Parameter sets: ML-DSA-44, ML-DSA-65, ML-DSA-87
% - Hints and compression: reducing signature size
% - Security proof sketch
% - Cite: FIPS 204, Ducas et al. (Dilithium paper)

\section{NTRU Lattices and the GPV Framework}
\label{sec:ntru-gpv}

% TODO: New section (foundation for Falcon)
% - NTRU lattices: definition and properties
% - The GPV framework: hash-and-sign from lattice trapdoors
% - Preimage sampling: Gaussian sampling with a trapdoor basis
% - The quality of the trapdoor: Gram-Schmidt norms
% - Cite: NTRU (Hoffstein et al. 1998), GPV 2008

\section{Falcon}
\label{sec:falcon}

% TODO: Expand from notes.tex §4.4.4
% - Construction: NTRU lattice + fast Fourier sampling
% - KeyGen: generating quality NTRU bases
% - Sign: tree-based Gaussian sampling (the "Falcon tree")
% - Verify: checking norm bounds
% - Parameters and security levels
% - Implementation challenge: constant-time floating-point
% - Cite: Fouque et al. (Falcon specification)

\section{ML-DSA vs Falcon: A Comparative Analysis}
\label{sec:sig-comparison}

% TODO: New section
% - Table: key sizes, signature sizes, speed (sign/verify)
% - Security assumptions: Module-LWE+SIS vs NTRU
% - Implementation complexity: Dilithium is simpler
% - Side-channel resistance: Dilithium is easier to protect
% - When to use which

%% Exercises
\begin{prob}
\label{prob:ch08-leakage}
Explain why a na\"ive lattice signature scheme that outputs $\mathbf{z} = \mathbf{s} \cdot c + \mathbf{y}$ (where $\mathbf{s}$ is the secret, $c$ the challenge, and $\mathbf{y}$ a random mask) leaks information about $\mathbf{s}$ after multiple signatures. How does rejection sampling solve this?
\end{prob}

\begin{prob}
\label{prob:ch08-dilithium}
In ML-DSA, the signer rejects and restarts if $\norm{\mathbf{z}}_\infty \geq \gamma_1 - \beta$. Explain the role of this check in security. What happens to the signing time as the bound $\gamma_1$ decreases?
\end{prob}
